\section{Decorrelation and the idempotent-product parent-commuting condition}%
\label{sec:decorrelation}

In \cite[Appendix~D.2]{Cirac2016MPDO_arXiv}, the parent commuting Hamiltonian
condition is stated for local projectors $Q_{AX}$ and $Q_{XB}$. With
$\widehat Q_{AX}$ and $\widehat Q_{XB}$ denoting their lifts to the tripartite space,
the commutation condition and common ground space $K_{AXB}$ satisfy
\begin{align}
    [\widehat Q_{AX},\widehat Q_{XB}] & = 0, \notag                                          \\
    K_{AXB}                           & = (K_{AX}\otimes H_B)\cap(H_A\otimes K_{XB}). \notag
\end{align}
Setting $P_{AX}=1-Q_{AX}$ and $P_{XB}=1-Q_{XB}$, the corresponding product
identity on the tripartite space is $\widehat P_{AX}\widehat P_{XB}=P_K$.
For a positive semidefinite operator on a bipartite space, the
support projector of the retained left marginal, lifted to the full space,
fixes the original operator, as used in
\cite[Appendix~D.2]{Cirac2016MPDO_arXiv}.
% Scope restriction: this section keeps the idempotent-product consequence
% and not the tensor-product locality or orthogonal-projector hypotheses of the
% source definition. See
% docs/paper-gaps/cpsv16_parent_commuting_hamiltonian_scope.tex.

\begin{lemma}[Partial trace of the identity]
    \label{lem:partial_trace_right_identity}
    \lean{Matrix.partialTraceRight_one}
    \leanok
    For finite-dimensional spaces $H_L$ and $H_R$,
    \[
        \tr_R(\mathbf 1_{L\otimes R})
        =(\dim H_R)\mathbf 1_L.
    \]
\end{lemma}

\begin{proof}
    \leanok
    In product-basis coordinates,
    \[
        \bigl(\tr_R\mathbf 1_{L\otimes R}\bigr)_{ij}
        =\sum_k\delta_{ij}
        =(\dim H_R)\delta_{ij}.
    \]
\end{proof}

\begin{lemma}[Adjoint identity for the right partial trace]
    \label{lem:trace_lift_left_factor_mul}
    \lean{Matrix.trace_leftKroneckerEmbed_mul}
    \leanok
    For every operator $M$ on $H_L$ and every operator $\rho$ on
    $H_L\otimes H_R$,
    \[
        \tr((M\otimes\mathbf 1_R)\rho)
        =\tr(M\tr_R\rho).
    \]
\end{lemma}

\begin{proof}
    \leanok
    Expand both traces in product bases and sum first over the identity entry
    on $H_R$.
\end{proof}

\begin{theorem}[Left absorption by the marginal support projector]
    \label{thm:marginal_support_left_absorption}
    \lean{Matrix.PosSemidef.leftKroneckerEmbed_supportProj_mul_self}
    \leanok
    \uses{lem:trace_lift_left_factor_mul}
    Let $\rho$ be positive semidefinite on $H_L\otimes H_R$, and let $P_L$ be
    the orthogonal projector onto the support of $\tr_R\rho$.
    Then
    \[
        (P_L\otimes\mathbf 1_R)\rho=\rho.
    \]
\end{theorem}

\begin{proof}
    \leanok
    Set $Q_L=\mathbf 1_L-P_L$.  Since
    $Q_L\tr_R\rho=0$, the adjoint identity gives
    \[
        \tr\bigl((Q_L\otimes\mathbf 1_R)\rho\bigr)
        =\tr(Q_L\tr_R\rho)=0.
    \]
    Positivity of $\rho$ then implies
    $(Q_L\otimes\mathbf 1_R)\rho=0$.  Finally,
    $\mathbf 1-(Q_L\otimes\mathbf 1_R)=P_L\otimes\mathbf 1_R$ gives the
    claim.
\end{proof}

\begin{theorem}[Right absorption by the marginal support projector]
    \label{thm:marginal_support_right_absorption}
    \lean{Matrix.PosSemidef.mul_leftKroneckerEmbed_supportProj_self}
    \leanok
    \uses{thm:marginal_support_left_absorption}
    Under the same assumptions,
    \[
        \rho(P_L\otimes\mathbf 1_R)=\rho.
    \]
    For $H_L=H_A\otimes H_X$ and $H_R=H_B$, the two absorption identities are
    precisely the $P_{AX}$ support-projector identities in
    \cite[Appendix~D.2, lines 2228--2235]{Cirac2016MPDO_arXiv}.
\end{theorem}

\begin{proof}
    \leanok
    Apply $(\cdot)^\dagger$ to the left absorption identity.  Since
    $\rho^\dagger=\rho$ and
    $(P_L\otimes\mathbf 1_R)^\dagger=P_L\otimes\mathbf 1_R$,
    \[
        \rho(P_L\otimes\mathbf 1_R)
        =\bigl[(P_L\otimes\mathbf 1_R)\rho\bigr]^\dagger
        =\rho^\dagger
        =\rho.
    \]
\end{proof}

\begin{definition}[Partial trace over the left factor]
    \label{def:partial_trace_left}
    \lean{Matrix.partialTraceLeft, Matrix.partialTraceLeft_apply}
    \leanok
    For an operator $\rho$ on $H_L\otimes H_R$, define
    \[
        (\tr_L\rho)_{r,s}
        =\sum_l \rho_{(l,r),(l,s)}.
    \]
\end{definition}

\begin{theorem}[Hermiticity of the left partial trace]
    \label{thm:partial_trace_left_hermitian}
    \lean{Matrix.partialTraceLeft_isHermitian}
    \leanok
    \uses{def:partial_trace_left}
    If $\rho=\rho^\dagger$, then
    $\tr_L\rho=(\tr_L\rho)^\dagger$.
\end{theorem}

\begin{proof}
    \leanok
    Entrywise Hermiticity gives
    \[
        \overline{(\tr_L\rho)_{r,s}}
        =\sum_k\overline{\rho_{(k,r),(k,s)}}
        =\sum_k\rho_{(k,s),(k,r)}
        =(\tr_L\rho)_{s,r}.
    \]
\end{proof}

\begin{theorem}[Positivity of the left partial trace]
    \label{thm:partial_trace_left_pos}
    \lean{Matrix.PosSemidef.partialTraceLeft}
    \leanok
    \uses{def:partial_trace_left}
    For finite-dimensional $H_L$ and $H_R$, if $\rho\geq 0$, then
    $\tr_L\rho\geq 0$.
\end{theorem}

\begin{proof}
    \leanok
    If $\rho^{(k)}$ is the principal submatrix indexed by
    $\{k\}\times H_R$, then
    \[
        (\tr_L\rho)_{r,s}
        =\sum_k\rho_{(k,r),(k,s)}
        =\sum_k\rho^{(k)}_{r,s}.
    \]
    Each $\rho^{(k)}$ is positive semidefinite, and hence so is their finite
    sum.
\end{proof}

\begin{lemma}[Adjoint identity for the left partial trace]
    \label{lem:trace_lift_right_factor_mul}
    \lean{Matrix.trace_rightKroneckerEmbed_mul}
    \leanok
    \uses{def:partial_trace_left}
    For every operator $M$ on $H_R$,
    \[
        \tr((\mathbf 1_L\otimes M)\rho)
        =\tr(M\tr_L\rho).
    \]
\end{lemma}

\begin{proof}
    \leanok
    Expanding in product bases gives
    \[
        \tr((\mathbf 1_L\otimes M)\rho)
        =\sum_{l,r}\sum_{k,s}\delta_{l,k}M_{rs}
          \rho_{(k,s),(l,r)}
        =\sum_{r,s}M_{rs}\sum_l\rho_{(l,s),(l,r)}
        =\tr(M\tr_L\rho).
    \]
\end{proof}

\begin{theorem}[Left absorption by the right marginal support]
    \label{thm:right_marginal_support_left_absorption}
    \lean{Matrix.PosSemidef.rightKroneckerEmbed_supportProj_mul_self}
    \leanok
    \uses{thm:partial_trace_left_pos,lem:trace_lift_right_factor_mul}
    Let $\rho$ be positive semidefinite.  If $P_R$ is the support projector of
    $\tr_L\rho$, then
    \[
        (\mathbf 1_L\otimes P_R)\rho=\rho.
    \]
\end{theorem}

\begin{proof}
    \leanok
    Set $Q=\mathbf 1_R-P_R$.  The preceding trace identity gives
    \[
        \tr\bigl((\mathbf 1_L\otimes Q)\rho\bigr)
        =\tr\bigl(Q\tr_L\rho\bigr)=0,
    \]
    since $P_R$ is the support projector of $\tr_L\rho$.
    Positivity of $\rho$ then gives
    $(\mathbf 1_L\otimes Q)\rho=0$, and hence
    $(\mathbf 1_L\otimes P_R)\rho=\rho$.
\end{proof}

\begin{theorem}[Right absorption by the right marginal support]
    \label{thm:right_marginal_support_right_absorption}
    \lean{Matrix.PosSemidef.mul_rightKroneckerEmbed_supportProj_self}
    \leanok
    \uses{thm:right_marginal_support_left_absorption}
    Under the same assumptions,
    \[
        \rho(\mathbf 1_L\otimes P_R)=\rho.
    \]
\end{theorem}

\begin{proof}
    \leanok
    Take adjoints in the left absorption identity.
\end{proof}

We now use the explicit tripartite indexing
$H_A\otimes(H_X\otimes H_B)$.  Reassociation identifies it with
$(H_A\otimes H_X)\otimes H_B$ whenever the $AX$ marginal is taken.

\begin{definition}[Tripartite local lifts]
    \label{def:tripartite_local_lifts}
    \lean{TripartiteDecorrelation.liftAX, TripartiteDecorrelation.liftXB,
      TripartiteDecorrelation.liftA, TripartiteDecorrelation.liftB}
    \leanok
    Writing $\widehat M_R$ for the lift of an operator $M_R$ supported on
    $R\in\{AX,XB,A,B\}$, the four lifts to
    $H_A\otimes(H_X\otimes H_B)$ are
    \begin{align*}
        \widehat M_{AX} &= M_{AX}\otimes\mathbf 1_B,
        & \widehat M_{XB} &= \mathbf 1_A\otimes M_{XB},\\
        \widehat M_A &= M_A\otimes\mathbf 1_{XB},
        & \widehat M_B &= \mathbf 1_{AX}\otimes M_B.
    \end{align*}
    The first and fourth formulas use the canonical identification
    $(H_A\otimes H_X)\otimes H_B\cong H_A\otimes(H_X\otimes H_B)$.
\end{definition}

\begin{lemma}[Products and adjoints of tripartite lifts]
    \label{lem:tripartite_local_lift_algebra}
    \lean{TripartiteDecorrelation.liftAX_mul,
      TripartiteDecorrelation.liftXB_mul,
      TripartiteDecorrelation.conjTranspose_liftAX,
      TripartiteDecorrelation.conjTranspose_liftXB,
      TripartiteDecorrelation.conjTranspose_liftA,
      TripartiteDecorrelation.conjTranspose_liftB}
    \leanok
    \uses{def:tripartite_local_lifts}
    The lifts from $AX$ and $XB$ preserve multiplication,
    \[
        \widehat{MN}_{AX}=\widehat M_{AX}\widehat N_{AX},
        \qquad
        \widehat{MN}_{XB}=\widehat M_{XB}\widehat N_{XB},
    \]
    and each of the four lifts from $AX$, $XB$, $A$, and $B$ preserves
    adjoints:
    \[
        (\widehat M_R)^\dagger=\widehat{M^\dagger}_R,
        \qquad R\in\{AX,XB,A,B\}.
    \]
\end{lemma}

\begin{proof}
    \leanok
    For the two-factor lifts,
    \[
        (M_{AX}\otimes\mathbf 1_B)(N_{AX}\otimes\mathbf 1_B)
        =(M_{AX}N_{AX})\otimes\mathbf 1_B,
        \qquad
        (\mathbf 1_A\otimes M_{XB})(\mathbf 1_A\otimes N_{XB})
        =\mathbf 1_A\otimes(M_{XB}N_{XB}).
    \]
    \[
        (M_{AX}\otimes\mathbf 1_B)^\dagger
          =M_{AX}^\dagger\otimes\mathbf 1_B,
        \qquad
        (\mathbf 1_A\otimes M_{XB})^\dagger
          =\mathbf 1_A\otimes M_{XB}^\dagger,
    \]
    and similarly
    \[
        (M_A\otimes\mathbf 1_{XB})^\dagger
          =M_A^\dagger\otimes\mathbf 1_{XB},
        \qquad
        (\mathbf 1_{AX}\otimes M_B)^\dagger
          =\mathbf 1_{AX}\otimes M_B^\dagger,
    \]
    with the tensor factors reassociated to the tripartite space.
\end{proof}

\begin{definition}[Tripartite observable decorrelation]
    \label{def:tripartite_observable_decorrelation}
    \lean{TripartiteDecorrelation.IsObservableDecorrelated}
    \leanok
    \uses{def:tripartite_local_lifts}
    Regions $A$ and $B$ are decorrelated relative to $P$ if, for all Hermitian
    observables $O_A=O_A^\dagger$ and $O_B=O_B^\dagger$,
    \[
        P O_A(\mathbf 1-P)O_B P=0.
    \]
    This is \cite[Definition~D.1, lines 2187--2191]{Cirac2016MPDO_arXiv}.
\end{definition}

\begin{definition}[Complex-linear extension of decorrelation]
    \label{def:tripartite_decorrelation}
    \lean{TripartiteDecorrelation.IsDecorrelated}
    \leanok
    \uses{def:tripartite_local_lifts}
    The complex-linear extension requires
    \[
        P O_A(\mathbf 1-P)O_B P=0
    \]
    for arbitrary complex matrices $O_A$ and $O_B$.
\end{definition}

\begin{theorem}[Observable and complex-linear decorrelation are equivalent]
    \label{thm:tripartite_observable_decorrelation_iff}
    \lean{TripartiteDecorrelation.isObservableDecorrelated_iff}
    \leanok
    \uses{def:tripartite_observable_decorrelation,
      def:tripartite_decorrelation}
    Observable decorrelation holds if and only if its complex-linear extension
    holds.
\end{theorem}

\begin{proof}
    \leanok
    Every complex matrix has the decomposition
    \[
        O=H+\mathrm{i}K,
        \qquad
        H=\frac{O+O^\dagger}{2},
        \qquad
        K=\frac{O-O^\dagger}{2\mathrm{i}},
    \]
    where $H=H^\dagger$ and $K=K^\dagger$.  For analogous decompositions
    $O=H+\mathrm{i}K$ and $O'=H'+\mathrm{i}K'$,
    \begin{align*}
        P O(\mathbf 1-P)O'P
        &=P H(\mathbf 1-P)H'P
          +\mathrm{i}P H(\mathbf 1-P)K'P\\
        &\quad+\mathrm{i}P K(\mathbf 1-P)H'P
          -P K(\mathbf 1-P)K'P.
    \end{align*}
    Each summand vanishes by observable decorrelation applied to its two
    Hermitian components.
\end{proof}

\begin{theorem}[Reverse order of decorrelation]
    \label{thm:tripartite_decorrelation_reverse}
    \lean{TripartiteDecorrelation.IsDecorrelated.reverse}
    \leanok
    \uses{def:tripartite_decorrelation}
    If $P=P^\dagger$ and $A$ and $B$ are decorrelated, then, for arbitrary
    complex matrices $O_A$ and $O_B$,
    \[
        P O_B(\mathbf 1-P)O_A P=0.
    \]
\end{theorem}

\begin{proof}
    \leanok
    Apply decorrelation to $O_A^\dagger$ and $O_B^\dagger$, and take adjoints.
\end{proof}

\begin{definition}[Tripartite reduced operators]
    \label{def:tripartite_marginals}
    \lean{TripartiteDecorrelation.marginalAX,
      TripartiteDecorrelation.marginalXB}
    \leanok
    The two reduced operators are
    \[
        \rho_{AX}=\tr_B P,
        \qquad
        \rho_{XB}=\tr_A P.
    \]
\end{definition}

\begin{theorem}[Matrix-unit expansions of the reduced operators]
    \label{thm:tripartite_marginal_matrix_units}
    \lean{TripartiteDecorrelation.lift_marginalAX_eq_sum,
      TripartiteDecorrelation.lift_marginalXB_eq_sum}
    \leanok
    \uses{def:tripartite_local_lifts,def:tripartite_marginals}
    For matrix units on $A$ and $B$,
    \[
        \mathbf 1_A\otimes\rho_{XB}
          =\sum_{a,k}E^A_{ak}P E^A_{ka},
        \qquad
        \rho_{AX}\otimes\mathbf 1_B
          =\sum_{b,k}E^B_{bk}P E^B_{kb}.
    \]
\end{theorem}

\begin{proof}
    \leanok
    For the $XB$ marginal, evaluation in the product basis gives
    \[
        \bigl(E^A_{ak}P E^A_{ka}\bigr)_{(a_0,r),(a_1,s)}
        =\delta_{a,a_0}\delta_{a,a_1}P_{(k,r),(k,s)}.
    \]
    Hence
    \[
        \sum_{a,k}\bigl(E^A_{ak}P E^A_{ka}\bigr)_{(a_0,r),(a_1,s)}
        =\delta_{a_0,a_1}\sum_kP_{(k,r),(k,s)}
        =\bigl(\mathbf 1_A\otimes\rho_{XB}\bigr)_{(a_0,r),(a_1,s)}.
    \]
    The $AX$ identity follows symmetrically.
\end{proof}

\begin{theorem}[Decorrelation for the reduced operators]
    \label{thm:tripartite_marginal_complementary_product_zero}
    \lean{TripartiteDecorrelation.IsDecorrelated.marginals_mul_complement_mul_eq_zero}
    \leanok
    \uses{def:tripartite_decorrelation,
      thm:tripartite_marginal_matrix_units}
    If $A,B$ are decorrelated, then
    \[
        \widehat\rho_{XB}(\mathbf 1-P)\widehat\rho_{AX}=0.
    \]
\end{theorem}

\begin{proof}
    \leanok
    Substitute the matrix-unit expansions.  Every summand has the form
    \[
        E^A_{ak}PE^A_{ka}(\mathbf 1-P)E^B_{bl}PE^B_{lb}
        =E^A_{ak}\bigl[PE^A_{ka}(\mathbf 1-P)E^B_{bl}P\bigr]E^B_{lb}
        =0,
    \]
    where the bracketed factor vanishes by decorrelation.
\end{proof}

\begin{theorem}[Reverse reduced-operator decorrelation]
    \label{thm:tripartite_reverse_marginal_complementary_product_zero}
    \lean{TripartiteDecorrelation.IsDecorrelated.reverse_marginals_mul_complement_mul_eq_zero}
    \leanok
    \uses{thm:tripartite_decorrelation_reverse,
      thm:tripartite_marginal_matrix_units}
    If $P=P^\dagger$ and $A,B$ are decorrelated, then
    \[
        \widehat\rho_{AX}(\mathbf 1-P)\widehat\rho_{XB}=0.
    \]
\end{theorem}

\begin{proof}
    \leanok
    Substitute the matrix-unit expansions.  Every summand has the form
    \[
        E^B_{bl}PE^B_{lb}(\mathbf 1-P)E^A_{ak}PE^A_{ka}
        =E^B_{bl}\bigl[PE^B_{lb}(\mathbf 1-P)E^A_{ak}P\bigr]E^A_{ka}
        =0,
    \]
    where the bracketed factor vanishes by reverse-order decorrelation.
\end{proof}

\begin{definition}[Tripartite marginal support projectors]
    \label{def:tripartite_marginal_supports}
    \lean{TripartiteDecorrelation.supportAX, TripartiteDecorrelation.supportXB}
    \leanok
    \uses{def:partial_trace_left}
    For a Hermitian operator $P$, let $P_{AX}$ and $P_{XB}$ be the support projectors of
    $\tr_B P$ and $\tr_A P$, respectively, as in
    \cite[Appendix~D.2, lines 2225--2229]{Cirac2016MPDO_arXiv}.  Write
    $\widehat P_{AX}$ and $\widehat P_{XB}$ for their lifts to the tripartite
    space.
\end{definition}

\begin{theorem}[Decorrelation for the two marginal supports]
    \label{thm:tripartite_complementary_product_zero}
    \lean{TripartiteDecorrelation.IsDecorrelated.supports_mul_complement_mul_eq_zero,
      TripartiteDecorrelation.IsDecorrelated.reverse_supports_mul_complement_mul_eq_zero}
    \leanok
    \uses{def:tripartite_marginal_supports,
      thm:tripartite_marginal_complementary_product_zero,
      thm:tripartite_reverse_marginal_complementary_product_zero}
    If $P=P^\dagger$ and $A,B$ are decorrelated, then
    \[
        \widehat P_{XB}(\mathbf 1-P)\widehat P_{AX}=0,
        \qquad
        \widehat P_{AX}(\mathbf 1-P)\widehat P_{XB}=0.
    \]
\end{theorem}

\begin{proof}
    \leanok
    For $g(0)=0$ and $g(x)=x^{-1}$ when $x\ne0$, functional calculus gives
    \[
        P_{AX}=g(\rho_{AX})\rho_{AX}=\rho_{AX}g(\rho_{AX}),
        \qquad
        P_{XB}=g(\rho_{XB})\rho_{XB}=\rho_{XB}g(\rho_{XB}).
    \]
    Hence
    \[
        \widehat P_{XB}(\mathbf 1-P)\widehat P_{AX}
        =\widehat{g(\rho_{XB})}
          \bigl[\widehat\rho_{XB}(\mathbf 1-P)\widehat\rho_{AX}\bigr]
          \widehat{g(\rho_{AX})}
        =0,
    \]
    and, in the reverse order,
    \[
        \widehat P_{AX}(\mathbf 1-P)\widehat P_{XB}
        =\widehat{g(\rho_{AX})}
          \bigl[\widehat\rho_{AX}(\mathbf 1-P)\widehat\rho_{XB}\bigr]
          \widehat{g(\rho_{XB})}
        =0.
    \]
    % Local fix for source lines 2236--2252: see
    % docs/paper-gaps/cpsv16_decorrelation_slice_expansion_fix.tex.
\end{proof}

\begin{theorem}[Absorption by the lifted $XB$ support projector]
    \label{thm:tripartite_xb_support_absorption}
    \lean{TripartiteDecorrelation.liftSupportXB_mul_self,
      TripartiteDecorrelation.mul_liftSupportXB_self}
    \leanok
    \uses{def:tripartite_local_lifts,def:tripartite_marginal_supports,
      thm:right_marginal_support_left_absorption,
      thm:right_marginal_support_right_absorption}
    For $P\geq0$,
    \[
        \widehat P_{XB}P=P=P\widehat P_{XB}.
    \]
\end{theorem}

\begin{proof}
    \leanok
    For the bipartition $A|XB$, the two support-absorption identities give
    \[
        (\mathbf 1_A\otimes P_{XB})P
        =P
        =P(\mathbf 1_A\otimes P_{XB}).
    \]
\end{proof}

\begin{theorem}[Absorption by the lifted $AX$ support projector]
    \label{thm:tripartite_ax_support_absorption}
    \lean{TripartiteDecorrelation.liftSupportAX_mul_self,
      TripartiteDecorrelation.mul_liftSupportAX_self}
    \leanok
    \uses{def:tripartite_local_lifts,def:tripartite_marginal_supports,
      thm:marginal_support_left_absorption,
      thm:marginal_support_right_absorption}
    For $P\geq0$,
    \[
        \widehat P_{AX}P=P=P\widehat P_{AX}.
    \]
\end{theorem}

\begin{proof}
    \leanok
    After reassociating the tripartite space as $AX|B$, the two
    support-absorption identities give
    \[
        (P_{AX}\otimes\mathbf 1_B)P
        =P
        =P(P_{AX}\otimes\mathbf 1_B).
    \]
\end{proof}

\begin{theorem}[Product of the two marginal support projectors]
    \label{thm:tripartite_support_products}
    \lean{TripartiteDecorrelation.supportProducts_eq}
    \leanok
    \uses{def:tripartite_marginal_supports,
      thm:tripartite_xb_support_absorption,
      thm:tripartite_ax_support_absorption,
      thm:tripartite_complementary_product_zero}
    If $P=P^\dagger=P^2$ and $A$ and $B$ are decorrelated relative to $P$, then
    \[
        \widehat P_{XB}\widehat P_{AX}
        =P
        =\widehat P_{AX}\widehat P_{XB}.
    \]
\end{theorem}

\begin{proof}
    \leanok
    Inserting $\mathbf 1=P+(\mathbf 1-P)$ between the two support projectors
    gives
    \[
        \widehat P_{XB}\widehat P_{AX}
        =\widehat P_{XB}P\widehat P_{AX}
          +\widehat P_{XB}(\mathbf 1-P)\widehat P_{AX}
        =P+0
        =P.
    \]
    Here the second equality uses the support-absorption identities and
    Theorem~\ref{thm:tripartite_complementary_product_zero}.  The same
    calculation with the two factors reversed gives
    $\widehat P_{AX}\widehat P_{XB}=P$.
\end{proof}

\begin{definition}[Parent commuting Hamiltonian on the tripartite space]
    \label{def:tripartite_parent_commuting_hamiltonian}
    \lean{TripartiteDecorrelation.HasGroundSpaceIntersection,
      TripartiteDecorrelation.CommutingParentHamiltonian,
      TripartiteDecorrelation.HasCommutingParentHamiltonian}
    \leanok
    \uses{def:tripartite_local_lifts}
    There are orthogonal projectors $P_{AX}$ and $P_{XB}$ on the two local
    factors whose lifts commute and whose ranges satisfy the source
    ground-space intersection condition
    \[
        \ran P
        =\ran\widehat P_{AX}
          \cap\ran\widehat P_{XB}.
    \]
    Thus $Q_{AX}=\mathbf 1-P_{AX}$ and
    $Q_{XB}=\mathbf 1-P_{XB}$ are exactly the commuting parent terms of
    \cite[Definition~D.2, lines 2205--2218]{Cirac2016MPDO_arXiv}.
\end{definition}

\begin{theorem}[Ground-space intersection and projector product]
    \label{thm:tripartite_ground_space_intersection_iff_product}
    \lean{TripartiteDecorrelation.range_mul_eq_inter,
      TripartiteDecorrelation.hermitian_idempotent_eq_of_range_eq,
      TripartiteDecorrelation.hasGroundSpaceIntersection_iff_product_eq,
      TripartiteDecorrelation.CommutingParentHamiltonian.hproduct}
    \leanok
    \uses{def:tripartite_parent_commuting_hamiltonian}
    Let $P$, $P_{AX}$, and $P_{XB}$ be orthogonal projectors, and write their
    lifts to the tripartite space as
    \[
        \widehat P_{AX}=P_{AX}\otimes\mathbf 1_B,
        \qquad
        \widehat P_{XB}=\mathbf 1_A\otimes P_{XB}.
    \]
    Suppose that the lifted projectors commute:
    \[
        \widehat P_{AX}\widehat P_{XB}
        =\widehat P_{XB}\widehat P_{AX}.
    \]
    Then
    \[
        \ran P
        =\ran\widehat P_{AX}\cap
          \ran\widehat P_{XB}
        \quad\Longleftrightarrow\quad
        \widehat P_{AX}\widehat P_{XB}=P.
    \]
    This is the passage from
    \cite[Definition~D.2, lines 2205--2218]{Cirac2016MPDO_arXiv} to the
    projector identity at lines 2279--2289.
\end{theorem}

\begin{proof}
    \leanok
    A vector lies in the range of
    $\widehat P_{AX}\widehat P_{XB}$ precisely when it is fixed by both
    commuting idempotents. Hence
    \[
        \ran(\widehat P_{AX}\widehat P_{XB})
        =\ran\widehat P_{AX}\cap
          \ran\widehat P_{XB}.
    \]
    The product is again an orthogonal projector.  Orthogonal projectors with
    equal ranges are equal, giving the forward implication; the reverse
    implication follows by substituting the product identity.
\end{proof}

\begin{theorem}[Parent commuting Hamiltonian--complex-linear decorrelation equivalence]
    \label{thm:tripartite_parent_iff_complex_linear_decorrelated}
    \lean{TripartiteDecorrelation.CommutingParentHamiltonian.isDecorrelated,
      TripartiteDecorrelation.IsDecorrelated.hasCommutingParentHamiltonian,
      TripartiteDecorrelation.parentHamiltonian_iff_decorrelated}
    \leanok
    \uses{def:tripartite_decorrelation,
      def:tripartite_parent_commuting_hamiltonian,
      thm:tripartite_support_products,
      thm:tripartite_ground_space_intersection_iff_product}
    Let $P=P^\dagger=P^2$.  Then $P$ is the ground-space projector of a
    parent commuting Hamiltonian on $AX$ and $XB$ if and only if regions
    $A$ and $B$ satisfy complex-linear decorrelation.
\end{theorem}

\begin{proof}
    \leanok
    Starting from decorrelation, take the support projectors of
    $\rho_{AX}$ and $\rho_{XB}$.  The preceding product theorem proves their
    commutation and identifies their product with $P$.

    Starting from commuting parent projectors, first use the ground-space
    intersection theorem to write
    $P=\widehat P_{AX}\widehat P_{XB}
      =\widehat P_{XB}\widehat P_{AX}$.  Locality gives
    $[O_A,\widehat P_{XB}]=[O_B,\widehat P_{AX}]=0$, while
    \[
        \widehat P_{XB}(\mathbf 1-P)\widehat P_{AX}=0.
    \]
    Therefore, using the two locality commutators,
    \[
    \begin{aligned}
        P O_A(\mathbf 1-P)O_B P
        &=(\widehat P_{AX}\widehat P_{XB})O_A(\mathbf 1-P)O_B
          (\widehat P_{AX}\widehat P_{XB})\\
        &=\widehat P_{AX}O_A
          \bigl[\widehat P_{XB}(\mathbf 1-P)\widehat P_{AX}\bigr]
          O_B\widehat P_{XB}\\
        &=0.
    \end{aligned}
    \]
\end{proof}

\begin{theorem}[Parent commuting Hamiltonian--observable decorrelation equivalence]
    \label{thm:tripartite_parent_iff_decorrelated}
    \lean{TripartiteDecorrelation.parentHamiltonian_iff_observableDecorrelated}
    \leanok
    \uses{def:tripartite_observable_decorrelation,
      thm:tripartite_observable_decorrelation_iff,
      thm:tripartite_parent_iff_complex_linear_decorrelated}
    Let $P=P^\dagger=P^2$.  Then $P$ is the ground-space projector of a
    parent commuting Hamiltonian on $AX$ and $XB$ if and only if regions
    $A$ and $B$ are decorrelated when tested on Hermitian observables.
    This is \cite[Proposition~D.3, lines 2221--2289]{Cirac2016MPDO_arXiv}.
\end{theorem}

\begin{proof}
    \leanok
    Combine the complex-linear parent equivalence with the decomposition
    $O=H+\mathrm{i}K$ from
    Theorem~\ref{thm:tripartite_observable_decorrelation_iff}.
\end{proof}

\begin{definition}[Decorrelation]\label{def:is_decorrelated}
    \lean{Decorrelation.IsDecorrelated}
    \leanok
    Given an endomorphism $P_K$ and sets of observables $\mathcal{O}_A$,
    $\mathcal{O}_B$, regions $A$ and $B$ are \emph{decorrelated} w.r.t.\ $K$
    when
    $P_K \circ O_A \circ (1 - P_K) \circ O_B \circ P_K = 0$
    for all $O_A \in \mathcal{O}_A$, $O_B \in \mathcal{O}_B$.
\end{definition}

\begin{definition}[Idempotent-product parent-commuting structure]\label{def:has_commuting_parent_ham}
    \lean{Decorrelation.HasCommutingParentHam}
    \leanok
    The structure used below records only the idempotent-product
    consequence of the parent commuting Hamiltonian condition of
    \cite[Appendix~D.2]{Cirac2016MPDO_arXiv}. It consists of commuting
    idempotent endomorphisms $P_{AX}$ and $P_{XB}$ such that, for the
    idempotent $P_K$ associated to $K$,
    \[
        P_{AX}\circ P_{XB}=P_K.
    \]
    In the source, this is the equation obtained from local projectors
    $Q_{AX}$ and $Q_{XB}$ satisfying $[Q_{AX},Q_{XB}]=0$ and
    $K_{AXB}=(K_{AX}\otimes H_B)\cap(H_A\otimes K_{XB})$.
    % Scope restriction: this definition records only the idempotent-product
    % part of the source condition. See
    % docs/paper-gaps/cpsv16_parent_commuting_hamiltonian_scope.tex.
\end{definition}

\begin{theorem}[Idempotence gives tautological idempotent-product data]%
    \label{thm:has_commuting_parent_ham_of_idem}
    \lean{Decorrelation.HasCommutingParentHam.ofIdem}
    \leanok
    \uses{def:has_commuting_parent_ham}
    If
    \[
        P_K \circ P_K = P_K,
    \]
    then $P_K$ has idempotent-product data.
\end{theorem}

\begin{proof}
    \leanok
    \uses{def:has_commuting_parent_ham}
    Take $P_{AX} = P_{XB} = P_K$. Their commutativity is immediate, and the
    product identity is exactly the assumed idempotence of $P_K$.
\end{proof}

\begin{lemma}[Cancellation for commuting idempotents]
    \label{lem:comp_complement_comm_zero}
    \lean{LinearMap.comp_complement_comm_zero}
    \leanok
    If $P$ and $Q$ are commuting idempotents, then
    \[
        Q \circ (1 - P \circ Q) \circ P = 0.
    \]
\end{lemma}

\begin{proof}
    \leanok
    The cancellation is the computation
    \[
        Q(1-PQ)P
        =
        QP-QPQP
        =
        QP-QPPQ
        =
        QP-QP
        =
        0,
    \]
    using commutativity and idempotence.
\end{proof}

\begin{theorem}[Idempotence of the product projector]%
    \label{thm:pk_idem}
    \lean{Decorrelation.HasCommutingParentHam.pK_idem}
    \leanok
    \uses{def:has_commuting_parent_ham}
    If $P_K = P_{AX} \circ P_{XB}$ is idempotent-product data, then
    \[
        P_K \circ P_K = P_K.
    \]
\end{theorem}

\begin{proof}
    \leanok
    \uses{def:has_commuting_parent_ham}
    Since $P_{AX}$ and $P_{XB}$ are idempotent and commute, their product
    $P_{AX} \circ P_{XB}$ is idempotent. Substituting
    $P_K = P_{AX} \circ P_{XB}$ gives the claim.
\end{proof}

\begin{theorem}[Absorption by the left projector]%
    \label{thm:pAX_comp_pK}
    \lean{Decorrelation.HasCommutingParentHam.pAX_comp_pK}
    \leanok
    \uses{def:has_commuting_parent_ham}
    If $P_K = P_{AX} \circ P_{XB}$ is idempotent-product data, then
    \[
        P_{AX} \circ P_K = P_K.
    \]
\end{theorem}

\begin{proof}
    \leanok
    \uses{def:has_commuting_parent_ham}
    Using $P_K = P_{AX} \circ P_{XB}$ and the idempotence of $P_{AX}$,
    \[
        P_{AX} \circ P_K
        = P_{AX} \circ P_{AX} \circ P_{XB}
        = P_{AX} \circ P_{XB}
        = P_K.
    \]
\end{proof}

\begin{theorem}[Absorption by the right projector]%
    \label{thm:pK_comp_pXB}
    \lean{Decorrelation.HasCommutingParentHam.pK_comp_pXB}
    \leanok
    \uses{def:has_commuting_parent_ham}
    If $P_K = P_{AX} \circ P_{XB}$ is idempotent-product data, then
    \[
        P_K \circ P_{XB} = P_K.
    \]
\end{theorem}

\begin{proof}
    \leanok
    \uses{def:has_commuting_parent_ham}
    Using $P_K = P_{AX} \circ P_{XB}$ and the idempotence of $P_{XB}$,
    \[
        P_K \circ P_{XB}
        = P_{AX} \circ P_{XB} \circ P_{XB}
        = P_{AX} \circ P_{XB}
        = P_K.
    \]
\end{proof}

\begin{theorem}[Cross-absorption by the right projector]%
    \label{thm:pXB_comp_pK}
    \lean{Decorrelation.HasCommutingParentHam.pXB_comp_pK}
    \leanok
    \uses{def:has_commuting_parent_ham}
    If $P_K = P_{AX} \circ P_{XB}$ is idempotent-product data, then
    \[
        P_{XB} \circ P_K = P_K.
    \]
\end{theorem}

\begin{proof}
    \leanok
    \uses{def:has_commuting_parent_ham}
    By commutativity,
    $P_{XB} \circ P_K = P_{XB} \circ P_{AX} \circ P_{XB}
        = P_{AX} \circ P_{XB} \circ P_{XB}$, and idempotence of $P_{XB}$
    yields $P_{XB} \circ P_K = P_K$.
\end{proof}

\begin{theorem}[Cross-absorption by the left projector]%
    \label{thm:pK_comp_pAX}
    \lean{Decorrelation.HasCommutingParentHam.pK_comp_pAX}
    \leanok
    \uses{def:has_commuting_parent_ham}
    If $P_K = P_{AX} \circ P_{XB}$ is idempotent-product data, then
    \[
        P_K \circ P_{AX} = P_K.
    \]
\end{theorem}

\begin{proof}
    \leanok
    \uses{def:has_commuting_parent_ham}
    By commutativity,
    $P_K \circ P_{AX} = P_{AX} \circ P_{XB} \circ P_{AX}
        = P_{AX} \circ P_{AX} \circ P_{XB}$, and idempotence of $P_{AX}$
    yields $P_K \circ P_{AX} = P_K$.
\end{proof}

\begin{theorem}[Reverse product identity]%
    \label{thm:reverse_product}
    \lean{Decorrelation.HasCommutingParentHam.reverse_product}
    \leanok
    \uses{def:has_commuting_parent_ham}
    If $P_K = P_{AX} \circ P_{XB}$ is idempotent-product data, then
    \[
        P_{XB} \circ P_{AX} = P_K.
    \]
\end{theorem}

\begin{proof}
    \leanok
    \uses{def:has_commuting_parent_ham}
    This follows from commutativity of $P_{AX}$ and $P_{XB}$ together
    with the identity $P_{AX} \circ P_{XB} = P_K$.
\end{proof}

\begin{theorem}[Commuting complementary projectors]%
    \label{thm:complement_comm}
    \lean{Decorrelation.HasCommutingParentHam.complement_comm}
    \leanok
    \uses{def:has_commuting_parent_ham}
    If $P_K = P_{AX} \circ P_{XB}$ is idempotent-product data and
    \[
        Q_{AX} = 1 - P_{AX}, \qquad Q_{XB} = 1 - P_{XB},
    \]
    then
    \[
        Q_{AX} \circ Q_{XB} = Q_{XB} \circ Q_{AX}.
    \]
\end{theorem}

\begin{proof}
    \leanok
    \uses{def:has_commuting_parent_ham}
    Expanding both products and using $P_{AX} \circ P_{XB}
        = P_{XB} \circ P_{AX}$ shows that the complementary projectors commute.
\end{proof}

\begin{theorem}[Idempotent-product data imply decorrelation]%
    \label{thm:commuting_ham_is_decorrelated}
    \lean{Decorrelation.commutingHam_isDecorrelated}
    \leanok
    \uses{def:is_decorrelated, def:has_commuting_parent_ham}
    If $P_{AX}$ and $P_{XB}$ are idempotent endomorphisms satisfying
    $P_{AX} \circ P_{XB} = P_K$ and
    $P_{AX} \circ P_{XB} = P_{XB} \circ P_{AX}$, and
    $A$-observables commute with $P_{XB}$ while $B$-observables commute
    with $P_{AX}$, then regions $A$ and $B$ are decorrelated w.r.t.\ $K$.
    This is the idempotent-product form of the backward direction of
    \cite[Appendix~D, Section~D.2, Proposition~D.3]{Cirac2016MPDO_arXiv}.
\end{theorem}

\begin{proof}
    \leanok
    \uses{lem:comp_complement_comm_zero}
    The key identity is $P_{XB} \circ (1 - P_{AX} \circ P_{XB}) \circ P_{AX} = 0$.
    Substituting $P_K = P_{AX} \circ P_{XB}$ and using
    $O_A \circ P_{XB} = P_{XB} \circ O_A$ and
    $O_B \circ P_{AX} = P_{AX} \circ O_B$ gives
    \[
        P_K \circ O_A \circ (1 - P_K) \circ O_B \circ P_K
        = P_{AX} \circ O_A \circ
          \bigl[P_{XB} \circ (1 - P_{AX} \circ P_{XB}) \circ P_{AX}\bigr]
          \circ O_B \circ P_{XB} = 0.
    \]
\end{proof}

\begin{theorem}[Decorrelating idempotent-product data]%
    \label{thm:has_commuting_parent_ham_is_decorrelated}
    \lean{Decorrelation.HasCommutingParentHam.isDecorrelated}
    \leanok
    \uses{thm:commuting_ham_is_decorrelated, def:has_commuting_parent_ham, def:is_decorrelated}
    Corollary of Theorem~\ref{thm:commuting_ham_is_decorrelated}: if
    $P_K$ has idempotent-product data, the $A$-observables
    commute with $P_{XB}$, and the $B$-observables commute with $P_{AX}$, then
    regions $A$ and $B$ are decorrelated w.r.t.\ $K$.
\end{theorem}

\begin{proof}
    \leanok
    \uses{thm:commuting_ham_is_decorrelated}
    Apply Theorem~\ref{thm:commuting_ham_is_decorrelated} to the witnessing
    endomorphisms $P_{AX}$ and $P_{XB}$ supplied by the idempotent-product data.
\end{proof}

\begin{lemma}[Frustration-free identity]\label{lem:frustration_free_ham_eq}
    \lean{LinearMap.frustration_free_ham_eq}
    \leanok
    For any endomorphisms $P$ and $Q$,
    \[
        (1-P) + (1-Q) - (1-P)\circ(1-Q) = 1 - P\circ Q.
    \]
\end{lemma}

\begin{proof}
    \leanok
    Expanding the left-hand side and cancelling gives the result.
\end{proof}

\begin{theorem}[Idempotent-product Hamiltonian identity]%
    \label{thm:hamiltonian_eq}
    \lean{Decorrelation.HasCommutingParentHam.hamiltonian_eq}
    \leanok
    \uses{def:has_commuting_parent_ham, lem:frustration_free_ham_eq}
    If $P_K = P_{AX} \circ P_{XB}$ is idempotent-product data, then
    \[
        Q_{AX} + Q_{XB} - Q_{AX} \circ Q_{XB} = 1 - P_K,
    \]
    where $Q_{AX} = 1 - P_{AX}$ and $Q_{XB} = 1 - P_{XB}$ are the
    complementary projectors.
    This is the algebraic identity obtained from the complementary projectors
    in \cite[Appendix~D, Section~D.2]{Cirac2016MPDO_arXiv}; it is not the full
    source parent commuting Hamiltonian condition.
\end{theorem}

\begin{proof}
    \leanok
    \uses{lem:frustration_free_ham_eq}
    Applying Lemma~\ref{lem:frustration_free_ham_eq} to $P_{AX}$ and
    $P_{XB}$, and substituting $P_{AX} \circ P_{XB} = P_K$.
\end{proof}

\begin{theorem}[Ground-space membership]\label{thm:mem_ground_iff}
    \lean{Decorrelation.HasCommutingParentHam.mem_ground_iff}
    \leanok
    \uses{def:has_commuting_parent_ham}
    If $P_K = P_{AX} \circ P_{XB}$ is idempotent-product data, then
    \[
        P_K v = v
        \quad\Longleftrightarrow\quad
        P_{AX} v = v \text{ and }  P_{XB} v = v.
    \]
    This is the algebraic consequence corresponding to equation~(D.2) from
    \cite{Cirac2016MPDO_arXiv}:
    $K_{AXB} = (K_{AX} \otimes H_B) \cap (H_A \otimes K_{XB})$.
    The source equation itself also carries the tensor-product local-subspace
    hypotheses.
\end{theorem}

\begin{proof}
    \leanok
    \uses{def:has_commuting_parent_ham, thm:pAX_comp_pK, thm:pXB_comp_pK}
    ($\Rightarrow$) If $P_K v = v$, then $P_{AX}(P_K v) = (P_{AX} \circ P_K) v
        = P_K v = v$ by left absorption, and similarly for $P_{XB}$.

    ($\Leftarrow$) If $P_{AX} v = v$ and $P_{XB} v = v$, then
    $P_K v = (P_{AX} \circ P_{XB}) v = P_{AX}(P_{XB} v) = P_{AX} v = v$.
\end{proof}

